% ============================================================
%  SECTION 6: EU HONEY MARKET AND REGULATORY IMPLICATIONS
% ============================================================

\section{EU Honey Market and Regulatory Implications}

The European honey market is the world's largest regional import
market, and its structural characteristics -- chronic domestic
production deficit, high dependence on imports from third countries,
and wide quality variation of supplied products -- constitute
a context in which regulatory decisions concerning honey
disqualification criteria have direct economic, trade, and consumer
consequences. This section analyzes this market context as
justification for the precise approach advocated in this article,
based on functional maturity and water activity, in place of an
overly simplified approach based on yeast count or morphological
criterion.

%,,,,,,,,,,,,,,,,,,,,,,,,,,,,,,
\subsection{Structure and dynamics of honey imports into the EU}
%,,,,,,,,,,,,,,,,,,,,,,,,,,,,,,

The European Union is the world's largest regional honey importer,
generating a trade balance deficit in the beekeeping sector that
has no equivalent in any other region~\cite{cbi2024,pollinatorhub2025}.
The own production of 27 member states amounts to an estimated
272,286 thousand tonnes per year (2022--2023 data), covering
only 55--60\% of domestic demand, estimated at approximately
480--500 thousand tonnes~\cite{pollinatorhub2025,ec_honey_autumn2024}.
The difference is filled by imports from third countries.

According to Eurostat data for 2023, the EU imported
\textbf{163,700~tonnes} of honey from non-EU countries at a value
of \textbf{EUR~359.3~million}; exports amounted to 24,900~tonnes
(EUR~146~million)~\cite{eurostat_honey2023}. Adding intra-Community
trade (honey moved between member states following prior import
from outside), the total volume of honey imported into the EU
supply chain is estimated at \textbf{300--360 thousand tonnes}
per year~\cite{cbi2024}.
In 2024, total imports (extra-EU + intra-Community) amounted to
approximately 290 thousand tonnes (a decline of 3.1\% year-on-year
after three years of growth), but the long-term trend remains
upward: the EU market is expected to reach a volume of 445 thousand
tonnes by 2035 (CAGR +1.9\%)~\cite{eurostat_honey2023}.

The main external suppliers in 2023 were:
\textbf{China} (37\% of extra-EU imports), \textbf{Ukraine} (28\%),
\textbf{Argentina} (approx.\ 7\%), \textbf{Mexico} and
\textbf{Cuba}~\cite{tridge2024}.
The largest national importers within the EU were Germany
(41 thousand tonnes, 25\% of extra-EU imports), Belgium
(31.4 thousand tonnes, 19\%), Poland (23.3 thousand tonnes, 14\%),
Spain (15.7 thousand tonnes), and France
(7.7 thousand tonnes)~\cite{eurostat_honey2023}.

%,,,,,,,,,,,,,,,,,,,,,,,,,,,,,,
\subsection{The scale of honey adulteration on the EU market: empirical data}
%,,,,,,,,,,,,,,,,,,,,,,,,,,,,,,

The scale of the honey authenticity problem on the EU market is
documented by a number of studies conducted at the institutional
level. The coordinated control action of the European Commission
``From the Hives'' (2021--2022), covering analysis of 320 samples
of imported honey, demonstrated that \textbf{46\%} of the analyzed
samples were suspected of non-compliance with the Honey Directive,
representing a threefold increase compared to a similar study from
2015--2017, when the figure was 14\%~\cite{eufromhives2021}.
The highest absolute percentage of suspect shipments concerned China
(74\%), while Turkey showed the highest relative rate
(93\%)~\cite{eufromhives2021}. Among the identified adulteration
practices, the addition of sugar syrups (from sugarcane, corn, rice)
dominated, along with blending authentic honey with syrup products
in proportions below the detection threshold of conventional
methods~\cite{eufromhives2021,schoder2026}.

These data indicate a real and serious honey authenticity problem
on the EU market, but simultaneously clearly profile the dominant
form of adulteration as \emph{sugar syrup adulteration}, not as
the distribution of honey with an elevated yeast count. JRC and OLAF
documents do not indicate a single case of honey disqualification
based solely on the yeast CFU/g count without simultaneous
determination of active fermentation or syrup
adulteration~\cite{eufromhives2021}. This is a significant
observation from the perspective of this article: adulteration
practice focuses on chemical dilution of honey, not on manipulation
of its microflora.

In response to this crisis, the European Commission adopted
Directive (EU) 2024/1438 (part of the ``Breakfast Directives''
package), introducing from 2026 an obligation to declare all
countries of origin of honey blend components in the principal
field of view of the label~\cite{cbi_traceability2024}. The
implementation of this requirement strengthens the traceability
infrastructure, but does not resolve the problem of the lack of
validated parameters specific to immaturity.

%,,,,,,,,,,,,,,,,,,,,,,,,,,,,,,
\subsection{Analytical gaps: absence of a validated immaturity biomarker}
%,,,,,,,,,,,,,,,,,,,,,,,,,,,,,,

Identification of immature honey as a distinct category of
adulteration requires the existence of analytical tools capable of
reliably distinguishing between immature and mature honey, and here
science encounters a fundamental limitation. As the review by
Tsagkaris et al.~\cite{tsagkaris2021} demonstrated, standard
physicochemical methods (moisture, diastase, invertase, HMF,
sugar profile) are not sufficient for unambiguous determination
of immaturity, because:

\begin{enumerate}
    \item \textbf{There is no specific marker:}
          no single chemical compound or enzyme is identified
          as a biomarker exclusively of immaturity in a manner
          independent of botanical and geographical type;
          the physicochemical parameters of mature honey from
          the tropics may resemble the parameters of immature
          honey from the temperate zone~\cite{tsagkaris2021,guo2020};
    \item \textbf{Metagenomics and metabolomics have limited
          discriminatory power without reference databases:}
          as Schoder~\cite{schoder2026} emphasizes in his 2026
          review, the high analytical resolution of omic methods
          (genomics, proteomics, NMR, GC-MS, LC-HRMS) does not
          guarantee unambiguous fraud attribution without solid,
          globally representative reference databases;
          \textit{``high analytical resolution alone therefore
          does not guarantee unambiguous fraud
          attribution''}~\cite{schoder2026};
    \item \textbf{Blending effects mask immaturity:}
          blending immature with mature honeys in proportions
          below 50\% is analytically indistinguishable from mature
          honey by conventional methods~\cite{schoder2026}; it
          should be emphasized that this problem also applies to
          the honey extraction stage -- it is impossible to
          determine what fraction of honey was mature and what
          was immature at the time of extraction;
    \item \textbf{Yeast presence is not a specific marker:}
          as demonstrated in Section~4, the osmophilic yeast count
          in honey is determined mainly by climatic production
          conditions and hive systems, not by the degree of
          maturity, which excludes CFU/g as a useful immaturity
          biomarker~\cite{rodriguezmachado2024,schoder2026}.
\end{enumerate}

The most recent studies point to promising directions in the search
for immaturity biomarkers: Sun et al.~\cite{sun2021} identified
decenedioic acid and other fatty acids of bee origin as potential
molecular markers whose concentration correlates with the degree of
maturity in a manner independent of yeast count. NMR-based
metabolomics~\cite{nmrhoney2023} enables simultaneous quantification
of several dozen compounds in a single measurement, providing
a basis for developing a multidimensional maturity index. However,
to date, none of the proposed biomarkers has undergone full
multilaboratory validation encompassing the global diversity of
botanical and geographical types~\cite{tsagkaris2021}.

It should be borne in mind that legally established methods should
be as simple, rapid, and accessible as possible, since they will
constitute the basic tool of every beekeeper (similar to the current
refractometer) for determining the timing of extraction, as it is
this stage that determines whether honey is mature or not.

%,,,,,,,,,,,,,,,,,,,,,,,,,,,,,,
\subsection{Regulatory risk: consequences of overly strict interpretation}
%,,,,,,,,,,,,,,,,,,,,,,,,,,,,,,

Adopting a disqualification criterion based on yeast CFU/g count,
without taking into account water activity, climatic production
conditions, and the absence of active fermentation, entails a number
of serious regulatory and market risks, whose analysis is necessary
for a proper assessment of the implications of the Apimondia
position:

\paragraph{Risk of discrimination against tropical producers.}
As demonstrated in Section~5, honeys from tropical and traditional
systems naturally contain higher concentrations of osmophilic yeasts
due to hive type, climatic conditions, and richer endemic
microbiological
biodiversity~\cite{tadele2025,stinglessyeastbrazil2021}.
Disqualification of these honeys based on CFU/g would effectively
result in excluding from the EU market products legally produced
in countries of East Africa, Southeast Asia, and Latin America,
while not affecting syrups-adulterated honeys (which contain a
standard or reduced yeast count following the ultrafiltration
process)~\cite{eufromhives2021}.

\paragraph{Risk of aggravating the deficit and price increases.}
Excluding from the EU market honeys from China, Ukraine, and
Argentina -- the three largest suppliers collectively accounting
for $\approx 72\%$ of extra-EU imports~\cite{tridge2024} -- based
on the CFU/g criterion would cause a sharp rise in honey prices
on the European market and an aggravation of the deficit,
particularly severe for consumers in countries with low domestic
consumption (Poland, Germany, Belgium as net
importers)~\cite{eurostat_honey2023}.

\paragraph{Risk of creating arbitrary trade barriers.}
Applying a criterion that has no basis in Codex Alimentarius or
EU Directive 2001/110/EC can be challenged by trading partners
under the WTO-SPS Agreement (Agreement on the Application of
Sanitary and Phytosanitary Measures), which requires that
trade-restrictive measures have a scientific basis and be
proportionate to the identified risk~\cite{codex_stan12,eudir2001110}.
The risk of harvesting immature honey exists for all honeys,
regardless of where the apiary is located. Introduced requirements
and restrictions should have a global dimension and their
proportional reflection in EU law with respect to EU honey producers.


\paragraph{Risk of missing the actual problem.}
As demonstrated by the ``From the Hives'' action, the dominant
form of honey adulteration in the EU is sugar syrup adulteration,
not the distribution of immature honey with an elevated yeast
count~\cite{eufromhives2021}.
Focusing inspection resources on CFU/g measurement at the expense
of isotope ratio analysis ($\delta^{13}$C, $\delta^{2}$H,
$\delta^{18}$O) and NMR fingerprinting would lead to allocation
of inspection resources inconsistent with the actual adulteration
risk profile.
It should be noted, however, that the European Commission later
withdrew from treating $^1$H-NMR profiling as a sufficiently
established official control method, emphasizing the need for
further research, validation, and harmonization of analytical
methods~\cite{EuropeanCommission2023NMR}.


%,,,,,,,,,,,,,,,,,,,,,,,,,,,,,,
\subsection{Regulatory proposals: harmonization, traceability, and proportionality}
%,,,,,,,,,,,,,,,,,,,,,,,,,,,,,,

Based on the above market and regulatory analysis, this article
formulates the following proposals addressed to standardization
bodies (ISO, Codex), EU institutions (DG SANTE, EC), and industry
organizations (Apimondia, Copa-Cogeca):

\begin{enumerate}
    \item \textbf{Update EU Directive 2001/110/EC with a water
          activity ($a_w$) criterion:}
          introduction of mandatory $a_w$ measurement alongside
          moisture content as a primary parameter for fermentation
          risk assessment. Measurement of $a_w$ is a validated,
          commercially available method (instruments such as
          AquaLab, Rotronic HygroLab) and inexpensive; this would
          constitute a substitution of the morphological criterion
          with a physicochemical criterion having direct scientific
          justification~\cite{ruegg1981,zamora2006};

    \item \textbf{Establish a multiparametric microbiological
          assessment protocol:} fermentation risk assessment should
          be conducted simultaneously for $a_w$, water content,
          storage temperature, and yeast count (in accordance with
          the risk equation (\ref{eq:fermentation_risk}) proposed
          in Section~4.3); none of these parameters may be
          decisive in isolation from the
          others~\cite{analytica2020,apinz2021};

%     \item \textbf{Implement and mandatorily apply traceability
%           systems throughout the supply chain:}
%           Directive (EU) 2024/1438 is a step in the right
%           direction, but insufficient; traceability down to the
%           apiary level with documentation of moisture and
%           temperature at each \textit{post-harvest} stage
%           is required (cold chain monitoring)~\cite{cbi_traceability2024,tsagkaris2021};

%     \item \textbf{Build global honey authenticity reference
%           databases:} as PMC~2026~\cite{schoder2026} indicates,
%           advanced omic methods are ineffective without
%           representative reference databases; the EU Reference
%           Laboratory for Honey (EURL) project should be extended
%           to include samples from tropical and traditional
%           systems;

%     \item \textbf{Regulate three product categories:}
%           (a)~functionally mature honey with natural yeast
%           microflora, free of restrictions; (b)~honey with
%           applied technological dehumidification, permitted
%           with process declaration; (c)~systemically immature
%           or adulterated honey, prohibited;
%           only this categorization is proportionate
%           to the actual risk profile on the
%           market~\cite{garcia2025,tadele2025,codex_stan12}.
\end{enumerate}
